\documentclass[fontset = fandol]{ctexart}

\ExplSyntaxOn \makeatletter

\keys_define:nn { hduthesis / exam }
  {
    watermark .tl_set:N = \l__hduthesis_exam_watermark_tl,
  }
\NewDocumentCommand \hdutexset {m} { \keys_set:nn { hduthesis / exam } {#1} }
\DeclareHookRule { shipout / background }
  { hduthesis / exam } { before } { pgfrcs }
\AddToHook { shipout / background } [ hduthesis / exam ]
  {
    \tl_if_empty:NF \l__hduthesis_exam_watermark_tl
      {
        \tikz [ remember~picture, overlay ]
          \node [ opacity = .2 ] at (current~page)
            {
              \includegraphics [ width = \linewidth ]
                { \l__hduthesis_exam_watermark_tl }
            };
      }
  }

\RequirePackage
  {
    geometry, mathtools, array, bm, esint, keytheorems, libertine, inconsolata,
    siunitx, fixdif, derivative, physics2, tasks, anyfontsize
  }
\RequirePackage[inline]{enumitem}
\RequirePackage{fontawesome5, fancyqr, circuitikz}
\hook_gput_code:nnn { cmd/faIcon/before } { . }
  {
    \group_begin:
    \fontsize { \fp_eval:n { .9 * \f@size } } { 0 } \selectfont
  }
\hook_gput_code:nnn { cmd/faIcon/after } { . } { \group_end: }
\geometry{margin = 1in}

\RequirePackage[svgnames]{xcolor}
\newkeytheorem{note}
  [ style = remark, numbered = false,
    name = \color{DarkRed}\sffamily \faIcon{info-circle}~Tip. ]
\newkeytheorem{problem}
  [ style = definition,
    name = \color{DarkGreen}\sffamily \faIcon{pen-square}~Problem ]
\newkeytheorem{solution}
  [ style = definition, numbered = false,
    name = \color{MidnightBlue}\sffamily \faIcon{check-square}~Solution ]
\usephysicsmodule{ab, op.legacy}
\usephysicsmodule{bm-um.legacy}
\settasks { label = \sffamily\Alph*., item-indent = 2.5em, label-offset = .5pt }
\newcommand \true {\addtocounter{task}{1}\faIcon{check-circle}}
\def \@result#1 {\underline{~#1~}}
\def \s@result {\textsf{【\qquad 】}}
\DeclareRobustCommand \result {\@ifstar\s@result\@result}
\keys_define:nn { textfig }
  {
    ratio      .tl_set:N    = \l__textfig_ratio_tl,
      ratio    .initial:n   = 7 \c_colon_str 2,
    position   .tl_set:N    = \l__textfig_position_tl,
      position .initial:n   = c,
    unknown    .code:n      = \tl_if_novalue:nF {#1}
      { \tl_set_eq:NN \l__textfig_position_tl \l_keys_key_tl }
  }
\NewDocumentEnvironment{textfig}{ O{} }
  {
    \group_begin:
    \keys_set:nn { textfig } {#1}
    \seq_set_split:Nee \l__whats_textfig_ratio_seq { \c_colon_str }
      { \l__textfig_ratio_tl }
    \noindent
    \minipage [ \l__textfig_position_tl ]
      {
        \fp_eval:n
          {
            .96*\seq_item:Nn \l__whats_textfig_ratio_seq {1}/
              (
                \seq_item:Nn \l__whats_textfig_ratio_seq {1} + 
                \seq_item:Nn \l__whats_textfig_ratio_seq {2}
              )
          }\dimexpr \linewidth
      }
  }
  {
    \endminipage
    \vspace{4.5pt plus 1.5pt minus 1.5pt} \par
    \group_end:
  }
\NewDocumentCommand \textfigsplit {}
  {
    \endminipage
    \hspace{.02\linewidth}
    \minipage [ \l__textfig_position_tl ]
      {
        \fp_eval:n
          {
            .96*\seq_item:Nn \l__whats_textfig_ratio_seq {2}/
              (
                \seq_item:Nn \l__whats_textfig_ratio_seq {1} + 
                \seq_item:Nn \l__whats_textfig_ratio_seq {2}
              )
          }\dimexpr \linewidth
      }
  }

% \hook_gput_code:nnn { cmd/maketitle/before } { . } { \vskip -\baselineskip }
% \hook_gput_code:nnn { cmd/maketitle/after } { . }
%   { \__weyoung_maketitle_aftercode_aux: }
\cs_new_protected:Nn \__weyoung_maketitle_aftercode_aux:
  {
    \vspace*{-2\baselineskip}
    \begin{center}
      \begin{minipage}{.66\textwidth}
        \centering
        \begin{tasks}
          [ label-width = 5\ccwd, item-indent = 2cm,
            after-item-skip = 0pt ](2)
          \task[\textbf{课程信息}：] 理学院/A0715022
          \task[\textbf{授课教师}：] 陈许敏
          \task[\textbf{解析提供}：] \textsf{Skyrmion}
          \task[{\makebox[5em][s]{\textbf{\textsf\TeX~排版}：}}] \textsf{Skyrmion}
        \end{tasks}
      \end{minipage}
      \begin{minipage}{.15\textwidth}
        \centering
        \begin{tikzpicture}
          \node [ inner~sep = 0 ] (wechat) {
            \fancyqr
              [ size = 4em, image~padding = 1,
                image = \raisebox{4em/33}{\textcolor{DarkGreen}{\faIcon{weixin}}},
                r~color = DarkGreen, version = 4
              ]{http://weixin.qq.com/r/sRJiemTEGjllrSAZ90ej}
            };
          \node [ inner~sep = .5ex, align = center, right, font = \scriptsize ]
            at (wechat.east) {未\\央\\学\\社};
        \end{tikzpicture}
      \end{minipage}
      \begin{minipage}{.15\textwidth}
        \centering
        \begin{tikzpicture}
          \node [ inner~sep = 0 ] (qq) {
            \fancyqr
              [ size = 4em, image~x~padding = 2, image~y~padding = 1,
                image = \raisebox{4em/33}{\textcolor{MidnightBlue}{\faIcon{qq}}},
                l~color = MidnightBlue, version = 4
              ]{https://qm.qq.com/q/PH7CsQWvm2}
            };
          \node [ inner~sep = .5ex, align = center, left, font = \scriptsize ]
            at (qq.west) {物\\理\\营};
        \end{tikzpicture}
      \end{minipage}
    \end{center}
  }

\newcommand \iu {{\mathup i\mkern1mu}}
\newcommand \eu {\mathup e}
\AtBeginDocument{%
  \let \Re \relax \let \Im \relax
  \DeclareMathOperator{\Re}{Re}
  \DeclareMathOperator{\Im}{Im}
}

\makeatother \ExplSyntaxOff
\usetikzlibrary{ fadings, patterns, arrows, decorations.markings }
\tikzset
  {
    > = stealth',
    every picture/.append style =
      { line join = round, line cap = round, thick }
  }

\renewcommand*\det[1]{|\mathbf{#1}|}
\newcommand*\mat[1]{\mathbf{#1}}
\newcommand*\inv[1]{#1^{-1}}
\newcommand*\trans[1]{#1^{\mkern-1.5mu\mathsf{T}}}
\newcolumntype{C}{@{}>{{}}c<{{}}@{}}

% \setCJKmainfont{LXGW WenKai}[AutoFakeSlant]
% \setCJKsansfont{LXGW Marker Gothic}[AutoFakeBold, AutoFakeSlant, Scale = .95]
% \setCJKmonofont{LXGW WenKai Mono}[AutoFakeBold, AutoFakeSlant]
\ctexset{section/format+ = \raggedright}
\usepackage{newtxmath}
% \usepackage[warnings-off={mathtools-colon,mathtools-overbracket}]{unicode-math}
% \setmathfont[AutoFakeBold]{STIX Two Math}

\makeatletter \def \@result #1{\underline{ \quad }} \makeatother
\DeclareMathOperator{\DFT}{DFT}

\title{\sffamily\bfseries 24Fall \textsc{Hdu} 「线性代数」补考试题}
\author{} \date{}

\begin{document}

\maketitle

\section{填空题（每题 3 分，共 15 分）}

\begin{problem}
  设 $\mat A =
  \begin{pmatrix}
    1 & 2 & 3\\
    2 & 4 & 6\\
    3 & 6 & 9
  \end{pmatrix}$，则 $|\mat A^*| =$ \result{$0$}.
\end{problem}
\begin{solution}

\end{solution}

\begin{problem}
  已知 $3$ 阶方阵 $\mat A$ 的特征值为 $1$, $2$, $3$，
  则 $|\mat A^3 - 5\mat A^2 + 7\mat A| =$ \result{$18$}.
\end{problem}
\begin{solution}
  
\end{solution}

\begin{problem}
  实二次型 $f(x_1,x_2,x_3) = x_1^2 + x_3^2 + 2x_1x_2 - 2x_2x_3$ 的秩为
  \result{$3$}.
\end{problem}
\begin{solution}
  
\end{solution}

\begin{problem}
  已知 $\rank(\mat A_{4 \times 3}) = 3$，
  $\mat B =
  \begin{pmatrix}
    2 & -1 & 3\\
    0 &  3 & 1\\
    0 &  0 & 0
  \end{pmatrix}$，则 $\rank(AB) =$ \result{$2$}.
\end{problem}
\begin{solution}
  
\end{solution}

\begin{problem}
  向量 $\bm\alpha = \trans{(1,2,3)}$
  与向量 $\bm\beta = \trans{(3,\lambda,1)}$ 正交，则 $\lambda =$ \result{$-3$}.
\end{problem}
\begin{solution}
  
\end{solution}

\section{选择题（每题 3 分，共 15 分）}

\begin{problem}
  设 $\mat A$ 为 $m$ 行 $n$ 列矩阵，$\mat B$ 为 $n$ 行 $m$ 列矩阵，则
  \begin{tasks}(2)
    \task 当 $m > n$ 时，必有 $|\mat{AB}| \neq 0$
    \task 当 $m > n$ 时，必有 $|\mat{AB}| = 0$
    \task 当 $m < n$ 时，必有 $|\mat{AB}| \neq 0$
    \task 当 $m < n$ 时，必有 $|\mat{AB}| = 0$
  \end{tasks}
\end{problem}
\begin{solution}
  
\end{solution}

\begin{problem}
  设 $\mat A$, $\mat B$ 都是 $n$ 阶对称矩阵，则下列结论 \emph{不正确} 的是
  \begin{tasks}(2)
    \task $\mat{AB}$ 也是对称矩阵
    \task $\mat A + \mat B$ 也是对称矩阵
    \task $\mat A^m + \mat B^m$（$m$ 为正整数）也是对称矩阵
    \task $\mat B \trans{\mat A} + \mat A \trans{\mat B}$ 也是对称矩阵.
  \end{tasks}
\end{problem}
\begin{solution}
  
\end{solution}

\begin{problem}
  二次型 $f(x_1,x_2,x_3) = x_1^2 + x_3^2 - 2x_1x_2 + 2x_1x_3 + 2x_2x_3$
  的正惯性指数为
  \begin{tasks}(4)
    \task $0$
    \task $1$
    \task $2$
    \task $3$
  \end{tasks}
\end{problem}
\begin{solution}
  
\end{solution}

\begin{problem}
  $n$ 阶方阵 $\mat A$ 可逆的\emph{充分必要}条件是
  \begin{tasks}(2)
    \task $\mat A\bm x = 0$ 有非零解
    \task $\mat A$ 中任意两个行向量都不成比例
    \task $\mat A$ 的每个行向量都是非零向量
    \task 对任何 $n$ 维非零向量 $\bm\alpha$，均有 $\mat A\bm\alpha \neq 0$
  \end{tasks}
\end{problem}
\begin{solution}
  
\end{solution}

\begin{problem}
  已知 $\eta_1$, $\eta_2$ 是非齐次线性方程组 $\mat A\bm x = \bm b$ 的两个不同的解，
  $\xi_1$, $\xi_2$ 是其导出组 $\mat A\bm x = 0$ 的一个基础解系，
  $k_1$, $k_2 \in \mathbb R$，则 $\mat A\bm x = \bm b$ 的通解为
  \begin{tasks}(2)
    \task $k_1\xi_1 + k_2(\xi_1 + \xi_2) + (\eta_2 - \eta_1)/2$
    \task $k_1\xi_1 + k_2(\xi_2 - \xi_1) + (\eta_2 + \eta_1)/2$
    \task $k_1\xi_1 + k_2(\eta_1 + \eta_2) + (\eta_2 - \eta_1)/2$
    \task $k_1\xi_1 + k_2(\eta_2 - \eta_1) + (\eta_2 + \eta_1)/2$
  \end{tasks}
\end{problem}
\begin{solution}
  
\end{solution}

\section{计算题（共 70 分）}

\begin{problem}
  已知向量组 $\bm\alpha_1 = \trans{(1,-1,1)}$, $\bm\alpha_2 = \trans{(1,-2,2)}$,
  $\bm\alpha_3 = \trans{(1,\lambda,5)}$ 为 $\mathbb R^3$ 空间中的一组基，
  求 $\lambda$.
\end{problem}
\begin{solution}
  
\end{solution}

\begin{problem}
  判断实二次型 $f(x_1,x_2,x_3) =
  \begin{pmatrix}
    x_1 & x_2 & x_3
  \end{pmatrix}
  \begin{pmatrix}
    3 & 1 & 0\\
    3 & 3 & 0\\
    0 & 0 & 1
  \end{pmatrix}
  \begin{pmatrix*}
    x_1\\ x_2\\ x_3
  \end{pmatrix*}$ 的正定性.
\end{problem}
\begin{solution}
  
\end{solution}

\begin{problem}
  设向量组 $\bm\alpha_1 = \trans{(1,1,2,-2)}$,
  $\bm\alpha_2 = \trans{(1,3,-\lambda,-2\lambda)}$,
  $\bm\alpha_3 = \trans{(1,-1,6,0)}$ 的秩为 $2$，求 $\lambda$ 的值.
\end{problem}
\begin{solution}
  
\end{solution}

\begin{problem}
  已知 $\mathbb R^3$ 的两组基分别为 $\bm\alpha_1 = \trans{(1,1,0)}$,
  $\bm\alpha_2 = \trans{(0,1,1)}$, $\bm\alpha_3 = \trans{(0,0,1)}$ 及
  $\bm\beta_1 = \trans{(1,-1,2)}$, $\bm\beta_2 = \trans{(1,1,-1)}$,
  $\bm\beta_3 = \trans{(-2,1,3)}$. 试求由基 $\bm\alpha_1$, $\bm\alpha_2$,
  $\bm\alpha_3$ 到基 $\bm\beta_1$, $\bm\beta_2$, $\bm\beta_3$ 的过渡矩阵 $\mat P$.
\end{problem}
\begin{solution}
  
\end{solution}

\begin{problem}
  设 $\bm\xi_1 = \trans{(1,1,1)}$, $\bm\xi_2 = \trans{(1,1,0)}$,
  $\bm\xi_3 = \trans{(1,-2,1)}$ 为 $3$ 阶方阵 $\mat A$ 的分别属于特征值
  $\lambda_1 = 0$, $\lambda_2 = 1$, $\lambda_3 = 3$ 的特征向量，试求方阵 $\mat A$.
\end{problem}
\begin{solution}
  
\end{solution}

\begin{problem}
  求向量组 $\bm\alpha_1 = \trans{(1,-1,0)}$, $\bm\alpha_2 = \trans{(2,1,3)}$, $\bm\alpha_3 = \trans{(3,1,0)}$, $\bm\alpha_4 = \trans{(5,0,7)}$ 的一个最大线性无关组，并将其余向量用该最大线性无关组线性表示.
\end{problem}
\begin{solution}
  
\end{solution}

\begin{problem}
  设向量组 $\bm\alpha_1 = \trans{(\lambda,2,10)}$,
  $\bm\alpha_2 = \trans{(-2,1,5)}$, $\bm\alpha_3 = \trans{(-1,1,4)}$,
  $\bm\beta = \trans{(1,1,2025)}$，问 $\lambda$ 为何值时，$\bm\beta$ 能由向量组
  $\bm\alpha_1$, $\bm\alpha_2$, $\bm\alpha_3$ 唯一线性表示？
\end{problem}
\begin{solution}
  
\end{solution}

\begin{problem}
  设矩阵 $\mat A =
  \begin{pmatrix}
    0 & 0 & 1\\
    x & 1 & y\\
    1 & 0 & 0
  \end{pmatrix}$
  有三个线性无关的特征向量，求 $x$, $y$ 应满足的条件.
\end{problem}
\begin{solution}
  
\end{solution}

\begin{problem}
  求一个正交变换 $\mat X = \mat{QY}$，把实二次型
  $f(x_1,x_2,x_3) = x_1^2 + 2x_2^2 + 3x_3^2 - 4x_1x_2 - 4x_2x_3$ 化为标准型.
\end{problem}
\begin{solution}
  
\end{solution}

\begin{problem}
  设 $\mat B$ 是一个三阶非零矩阵，它的每一列都是齐次线性方程组
  \[
    \begin{cases}
      \begin{array}{@{}rCrCrCr@{}}
        x_1  & + & 2x_2 & - & 2x_3    & = & 0\\
        2x_1 & - & x_2  & + & \lambda & = & 0\\
        3x_1 & + & x_2  & - & x_3     & = & 0
      \end{array}
    \end{cases}
  \]
  的解，求 $\lambda$ 的值和 $|\mat B|$.
\end{problem}
\begin{solution}
  
\end{solution}

\end{document}